% !TEX root = ../main.tex
\begin{hyphenrules}{spanish}
La mayoría de los principales modelos de lenguaje abiertos de 2026 usan la arquitectura de mezcla de expertos (MoE)\@. En una capa MoE, un enrutador envía cada token solo a unos pocos de los muchos expertos. Muchos de estos modelos equilibran la carga de los expertos con un sesgo por experto, que baja si el experto recibe demasiados tokens y sube si recibe demasiado pocos. Además, la mayoría ha sustituido el optimizador AdamW por Muon, que ortogonaliza la actualización de cada matriz de pesos. Sin embargo, los laboratorios que los desarrollan discrepan sobre si Muon debe entrenar también el enrutador. Tampoco hay estudios publicados que midan cómo afecta el optimizador al esfuerzo del sesgo. Medimos este esfuerzo con la amplitud, la diferencia entre el mayor y el menor sesgo de una capa. Entrenamos modelos MoE pequeños, de 55 y 108 millones de parámetros activos, con texto web educativo. En 145 entrenamientos, variamos el optimizador de cada parte de la red. Salvo en el ajuste de las tasas de aprendizaje, entrenamos cada configuración con tres semillas aleatorias.

Con Muon, el sesgo se esforzó más que con AdamW\@. La amplitud fue mayor en todas las semillas y más del doble en la primera capa MoE\@. Con la mitad de la tasa de aprendizaje de Muon, desaparecieron unas tres cuartas partes de la amplitud adicional de esta capa. Muon logró además una pérdida menor y expertos más especializados. La amplitud de la primera capa MoE dependió del optimizador del resto de la red, no del optimizador del enrutador. Con el resto de la red en Muon, las entradas del primer enrutador compartían una dirección, con la que este favorecía a los mismos expertos en cada token. El sesgo tardó cientos de pasos en compensar esta preferencia. Mientras tanto, el experto más ocupado recibió de cuatro a siete veces más tokens de los que le correspondían. Al quitar esta dirección de la entrada del primer enrutador en los primeros 3000 pasos, la sobrecarga y la mayor parte de la amplitud adicional desaparecieron, y la pérdida no cambió. Con el enrutador en AdamW, las cargas variaron más al inicio del entrenamiento que con Muon, y la pérdida fue algo mayor. Ningún entrenamiento divergió, ni siquiera multiplicando por 16 las tasas de aprendizaje ajustadas. La amplitud siguió siendo mayor con un sesgo más rápido, al añadir una pérdida de equilibrado y en el modelo grande.

Para modelos de nuestro tamaño, estos resultados aconsejan mantener el enrutador en Muon si el resto de la red está en Muon. Además, señalan que conviene vigilar sobre todo la primera capa MoE.
\end{hyphenrules}
